\section{Method}
\label{sec:method}

\subsection{Architecture}

Figure~\ref{fig:pipeline} shows the pipeline.
A YAML configuration drives all parameters.
The reference genome is loaded and shared via \texttt{Arc}.
Variants come from VCF or stochastic generation.
The genome is partitioned into regions (one per chromosome by default), each processed by a Rayon\autocite{rayon} worker thread.
Workers send completed batches through a bounded \texttt{crossbeam-channel} (default depth 64) to a dedicated writer thread that produces FASTQ, BAM, and truth VCF output.

\begin{figure}[t]
\centering
\begin{tikzpicture}[
    node distance = 1.2cm and 1.8cm,
    box/.style    = {draw, rounded corners=3pt, minimum width=2.6cm,
                     minimum height=0.7cm, align=center, font=\small},
    io/.style     = {draw, trapezium, trapezium left angle=70,
                     trapezium right angle=110, minimum width=2.2cm,
                     minimum height=0.6cm, align=center, font=\small},
    worker/.style = {draw, dashed, rounded corners=5pt, inner sep=6pt},
    arr/.style    = {-Stealth, thick},
  ]
  \node[io, fill=blue!10] (cfg) {YAML Config};
  \node[io, fill=blue!10, right=1.5cm of cfg] (ref) {Reference};
  \node[io, fill=blue!10, right=1.5cm of ref] (vcf) {VCF};
  \node[box, fill=blue!5, below=1.0cm of ref] (parse) {Config + Reference\\Loading};
  \draw[arr] (cfg) -- (parse);
  \draw[arr] (ref) -- (parse);
  \draw[arr] (vcf) -- (parse);
  \node[box, fill=orange!15, below=0.9cm of parse] (part) {Region Partitioning};
  \draw[arr] (parse) -- (part);
  \node[box, fill=green!10, below left=1.0cm and -0.3cm of part, minimum width=1.9cm] (frag) {Fragment\\Sampling};
  \node[box, fill=green!10, right=0.4cm of frag, minimum width=1.9cm] (spike) {Variant\\Spike-In};
  \node[box, fill=green!10, right=0.4cm of spike, minimum width=1.9cm] (qual) {Quality\\Model};
  \node[box, fill=green!10, right=0.4cm of qual, minimum width=1.9cm] (art) {UMI \&\\Artefacts};
  \begin{scope}[on background layer]
    \node[worker, fill=green!3, fit=(frag)(spike)(qual)(art),
          label={[font=\footnotesize\itshape]above:Rayon workers (one per region)}] {};
  \end{scope}
  \draw[arr] (part) -- (frag);
  \draw[arr] (frag) -- (spike);
  \draw[arr] (spike) -- (qual);
  \draw[arr] (qual) -- (art);
  \node[box, fill=orange!25, below=1.2cm of spike, minimum width=3.0cm] (chan) {Bounded Channel\\(depth=64)};
  \draw[arr] (art) -- (chan);
  \node[box, fill=red!10, below=0.8cm of chan, minimum width=2.5cm] (writer) {Writer Thread};
  \draw[arr] (chan) -- (writer);
  \node[io, fill=yellow!20, below left=1.0cm and 1.2cm of writer] (fq) {FASTQ.gz};
  \node[io, fill=yellow!20, below=1.0cm of writer] (bam) {BAM};
  \node[io, fill=yellow!20, below right=1.0cm and 1.2cm of writer] (truth) {Truth VCF};
  \draw[arr] (writer) -- (fq);
  \draw[arr] (writer) -- (bam);
  \draw[arr] (writer) -- (truth);
\end{tikzpicture}
\caption{VarForge processing pipeline. Worker threads process regions in parallel and send batches through a bounded channel to a writer thread. Memory is bounded by channel depth, not dataset size.}
\label{fig:pipeline}
\end{figure}

\subsection{Read generation}

Fragment lengths follow $\mathcal{N}(\mu, \sigma^2)$ for standard libraries or a cfDNA mixture model (Section~\ref{sec:cfdna}).
Start positions are drawn uniformly within each region.
Two quality models are supported: a parametric model with per-cycle decay ($\bar{Q}(c) = Q_0 - \lambda c$), and an empirical model learned from a real BAM via \texttt{learn-profile}.

\subsection{Stochastic VAF sampling}

For each read overlapping a variant site, \varforge{} draws independently:
\begin{equation}
  X_i \sim \text{Bernoulli}(\text{VAF}_\text{eff})
  \label{eq:bernoulli}
\end{equation}
The alt-read count follows $A \sim \text{Binomial}(D, \text{VAF}_\text{eff})$, matching the natural sampling variance of real sequencing.
BAMSurgeon uses deterministic $A = \lfloor D \cdot \text{VAF} \rfloor$.

\subsection{Variant spike-in}

SNVs replace reference bases.
Indels modify sequence and CIGAR with local re-alignment.
MNVs are applied atomically to preserve cis-phasing.
Structural variants (deletions, insertions, inversions, tandem duplications, translocations) produce split reads and discordant pairs with correct CIGAR strings, SA tags, and mate coordinates.
Copy number variants scale read depth by $(p \cdot \text{CN}_\text{tumor} + (1-p) \cdot 2) / 2$ and adjust effective VAF accordingly.

\subsection{Tumour modelling}

Heterogeneity is a directed tree of clones, each with a cancer cell fraction (CCF).
The effective VAF integrates purity, ploidy, CCF, and copy number:
\begin{equation}
  \text{VAF}_\text{eff} = \frac{\text{CCF} \times m \times p}{p \cdot \text{CN}_\text{tumor} + (1-p) \cdot \text{CN}_\text{normal}}
  \label{eq:vaf}
\end{equation}

\subsection{UMI and duplex simulation}

Each molecule receives a UMI barcode of configurable length.
Family sizes follow a log-normal distribution reflecting real PCR kinetics.
In duplex mode\autocite{schmitt2012duplex}, each double-stranded molecule yields two strand families with swapped UMI pairs $(\alpha,\beta)$ and $(\beta,\alpha)$.
True mutations appear on both strands.
Artefacts are strand-specific, reproducing the error suppression that makes duplex sequencing valuable.

\subsection{Cell-free DNA fragmentation}
\label{sec:cfdna}

Fragment lengths follow a four-component Gaussian mixture:
\begin{equation}
  P(L) = \sum_{k=1}^{4} w_k \cdot \mathcal{N}(L;\, \mu_k,\, \sigma_k^2)
  \label{eq:cfdna}
\end{equation}
Components represent sub-nucleosomal (${\sim}143$\,bp), mono-nucleosomal (${\sim}167$\,bp), di-nucleosomal (${\sim}334$\,bp), and tri-nucleosomal (${\sim}500$\,bp) populations.
Below 167\,bp, a 10\,bp sinusoidal periodicity matching the DNA helical repeat is applied\autocite{snyder2016cfdna_nucleosome}.
Tumour-derived fragments are shifted shorter (${\sim}143$--147\,bp peak)\autocite{lo2010cfdna}, enabling fragment-based tumour fraction estimation by tools such as DELFI.

\subsection{Library artefacts}

FFPE deamination: C$>$T transitions at a configurable rate, preferentially at fragment ends.
Oxidative damage: G$>$T transversions with strand-specific bias.
PCR errors: per-cycle substitution errors propagated through family genealogies.
GC bias: fragment-level coverage modulation with an empirical U-shaped weighting.

\subsection{Implementation}

\varforge{} uses a pure-Rust dependency stack: noodles\autocite{noodles} for BAM/VCF/FASTA, \texttt{seq\_io} for FASTQ, \texttt{flate2} for compression.
No C libraries are required, enabling single-binary distribution.
Deterministic seeding uses $\text{seed}_r = H(\text{seed}_\text{global}, r)$ per region, ensuring reproducibility independent of thread scheduling.
